\documentclass[12pt,letterpaper]{article}

% ---- Packages ----
\usepackage[utf8]{inputenc}
\usepackage[T1]{fontenc}
\usepackage{amsmath,amssymb}
\usepackage{newtxtext,newtxmath}  % Times New Roman text + math (replaces mathptmx)
\usepackage[margin=1in]{geometry}
\usepackage{setspace}
\usepackage{titlesec}
\usepackage{enumitem}
\usepackage{xr-hyper}
\usepackage{hyperref}
\externaldocument{3 - Gravity}
\externaldocument{7 - Quantum Mechanics}
\usepackage{microtype}
\usepackage{booktabs}

% ---- Section formatting ----
\titleformat{\section}{\large\bfseries\sffamily}{\thesection.}{0.5em}{}
\titleformat{\subsection}{\normalsize\bfseries\sffamily}{\thesubsection}{0.5em}{}
\titleformat{\subsubsection}{\normalsize\bfseries}{\thesubsubsection}{0.5em}{}

% ---- Spacing ----
\setstretch{1.15}
\setlength{\parskip}{0.5em}
\setlength{\parindent}{0em}

% ---- Custom commands ----
\newcommand{\rhoa}{\rho_{\!A}}
\newcommand{\Ka}{K_{\!A}}
\newcommand{\vin}{v_{\mathrm{in}}}
\newcommand{\rs}{r_s}
% \hbar is already defined in LaTeX; no redefinition needed

\begin{document}

\begin{center}
{\LARGE\bfseries Unifying Theory of Dynamic Aether Mechanics:\\
Special Relativity}\\[1em]
{\large Erik Otto}
\end{center}

\vspace{1em}

\begin{abstract}
Within the Dynamic Aether Mechanics framework, special relativity emerges
naturally from the properties of a compressible quantum superfluid medium.  This
paper derives the Lorentz factor $\gamma = 1/\sqrt{1 - v^2/c^2}$ from geometric
path-length effects in the medium, time dilation from both velocity (path geometry)
and gravity (inflow-velocity-dependent time dilation), and mass--energy equivalence $E = mc^2$
from the compression wave energy of dissolving mass.  The five classical experiments
historically regarded as having disproved aether theories---Michelson--Morley,
Kennedy--Thorndike, Sagnac, stellar aberration, and Fizeau---are all quantitatively
predicted by the model.  The linearity of the medium's pressure--density response is
confirmed to $10^{-8}$ precision across velocities from $10^{-6}\,c$ to $0.9994c$.
\end{abstract}

\vspace{1em}

\bigskip
\noindent\textit{This paper is part of a series presenting the Unifying Theory of
Dynamic Aether Mechanics.  For the foundational concepts of Aether binding,
the three independent properties (density, pressure, temperature), and the
unification of force constants, see Paper~1: General Overview \cite{otto-overview}.}
\bigskip

\section{Special Relativity}
\label{sec:special-relativity}
% ============================================================

The greatest challenge to this theory is not internal failures, but the extraordinary
success of special relativity.\ However, many effects predicted by special relativity also
emerge naturally from dynamic Aether.


\subsection{Length Contraction}

Every force that binds matter---electromagnetic, strong, weak---propagates
through the Aether at the local speed of sound $c = \sqrt{\Ka/\rhoa}$.\ The
wave equation governing this propagation,
$\partial^2\phi/\partial t^2 = c^2\,\nabla^2\phi$, is invariant under the
Lorentz transformation \cite{lorentz1904}.\ This invariance guarantees that
if a bound system is in equilibrium at rest, the transformed
configuration---contracted by $1/\gamma$ in the direction of motion, with
perpendicular dimensions unchanged---is also a valid equilibrium for the same
system moving at $v$ through the medium.\ Lorentz \cite{lorentz1904} proved
this rigorously for electromagnetic forces; in this theory the result extends
to all matter, because every binding force propagates through the Aether
at~$c$.\ Length contraction is therefore not an \textit{ad hoc}
hypothesis but a mathematical consequence of all forces sharing the same
medium and the same propagation speed.

The anisotropy that drives the contraction is visible in the force propagation
times of a moving system.\ Force signals traveling along the direction of motion
encounter asymmetric propagation: upstream at effective speed $c - v$,
downstream at $c + v$.\ The round-trip signal time along a separation $L$
parallel to the motion is
%
\begin{equation}
  t_{\parallel} = \frac{L}{c - v} + \frac{L}{c + v}
    = \frac{2L/c}{1 - v^2/c^2}
\end{equation}
%
while the round-trip time across the same separation perpendicular to the
motion is
%
\begin{equation}
  t_{\perp} = \frac{2L}{\sqrt{c^2 - v^2}}
    = \frac{2L/c}{\sqrt{1 - v^2/c^2}}
\end{equation}
%
These differ by a factor of $\gamma$: force signals along the direction of
motion are delayed relative to perpendicular signals.\ In the contracted
equilibrium ($L' = L/\gamma$ parallel, $L$ perpendicular), this anisotropy is
absorbed into the steady-state field configuration---the retarded fields of
all the moving constituents self-consistently produce the same binding as in
the rest frame.

Note that this contraction is experimentally confirmed only for
\emph{velocity-driven} motion through the Aether.\ In the gravitational case,
the Aether's natural coordinate system (Painlev\'{e}--Gullstrand coordinates)
has flat spatial slices: the inflow produces time dilation
(Section~\ref{sec:grav-time-dilation}) but no radial spatial distortion.\
Gravitational length contraction has no experimental confirmation.

A secondary consequence of velocity contraction is compression of the Aether
in the rest frame of the moving object.\ Because the object does real work on
the medium, a genuine density increase $\delta\rho/\rhoa \approx v^2/(2c^2)$
develops via Bernoulli (see Paper~7~\cite{otto-quantum}, Section~6.2 for the
vortex-level treatment).\ This compression is a consequence of the
contraction, not its cause---the cause is the Lorentz invariance of the wave
equation.\ The compression is physically distinct from gravitational inflow,
where free-fall kinetic energy exactly balances potential energy, leaving
$\Delta\rho = 0$.


\subsection{Time Dilation---Velocity}

All internal processes of an object---atomic oscillations, wave propagation within its
structure, binding cycles---depend on waves that propagate at speed $c$ relative to the
Aether, not relative to the object.\ When the object moves at velocity $v$ through the
Aether, these internal waves must travel longer geometric paths through the medium.

Consider any wave-based process within the moving object.\ A wave crossing the object
perpendicular to its direction of motion covers distance $ct$ in the Aether frame while the
object moves $vt$.\ By the Pythagorean theorem, the perpendicular distance covered by the
wave is only $\sqrt{c^2 - v^2}\,t$.\ The internal process that took time $t_0$ at rest now
takes time $t_0 / \sqrt{1 - v^2/c^2}$---exactly the Lorentz factor $\gamma$.

Since every force in this theory (electric, magnetic, gravitational) propagates through
Aether at speed $c$, every internal process of the moving object is slowed by the same
factor.\ Time is fundamentally a measure of relative particle movement (a year measures
Earth's orbital motion; a clock measures the motion of its internal mechanisms), so this
universal slowing produces the observed time dilation.

Note that time dilation and length contraction share the same root cause---motion
through a medium with finite propagation speed $c$---and both follow from the
Lorentz invariance of the wave equation governing that medium.\ Time dilation
arises from longer \emph{wave paths} (Pythagorean theorem); length contraction
arises from the \emph{anisotropic retarded-field configuration} of a moving
bound system, whose equilibrium contracts by $1/\gamma$ in the direction of
motion \cite{lorentz1904}.

The factor $v^2/c^2$ that governs time dilation is the same factor that appears in the
magnetic-to-electric force ratio ($F_{\mathrm{mag}} = F_{\mathrm{elec}} \cdot v_1 v_2/c^2$) and in relativistic energy
($E = \gamma mc^2$).\ All four arise from the same Aether mechanics: the ratio of motion
velocity to medium propagation speed.\ This connection is developed in the Unification of
Force Constants section.


\subsection{Time Dilation---Gravitational}
\label{sec:grav-time-dilation}

Near massive objects, the Aether flows inward---drawn by the collective electron
vortices of the mass through the transient bind--unbind cycle described in
Section~\ref{sec:gravity}.\ This inflow is the direct cause of gravitational time
dilation.\ A static observer---one who maintains a fixed position while the Aether
flows past---has all internal processes affected by the relative motion between
their wave-based structure and the medium.\ Because internal processes depend on
waves propagating at the local speed $c = \sqrt{\Ka/\rhoa}$ \emph{relative to the
Aether}, a static observer in a flowing medium experiences a slowing of all
processes identical to velocity time dilation.

The inflow velocity at distance $r$ from a mass $M$ follows from energy
conservation (an Aether element falling from rest at infinity):
%
\begin{equation}
  \vin(r) = \sqrt{\frac{2GM}{r}}
  \label{eq:inflow}
\end{equation}
%
A static observer at radius $r$ is stationary while the Aether flows past at
$\vin$.\ The time dilation experienced by this observer is:\footnote{A
density-based approach would require $\Delta\rho/\rhoa = 2GM/(rc^2)$ to
reproduce the observed time dilation from $c$-variation alone.\ However,
Bernoulli's equation for the Aether's logarithmic equation of state shows
this is unphysical: with inflow at $\vin = \sqrt{2GM/r}$, the kinetic energy
exactly balances the gravitational potential energy, leaving zero energy for
compression ($\Delta\rho/\rhoa = 0$ exactly).\ The maximum density perturbation
(hydrostatic equilibrium, $v = 0$) is $GM/(rc^2)$---half the required value.\
The flow-based derivation resolves this by attributing time dilation to the
inflow velocity rather than to a density change.}
%
\begin{equation}
  \frac{\Delta t}{t} = \frac{\vin^2}{2\,c^2} = \frac{GM}{r\,c^2}
\end{equation}
%
This matches the observed gravitational time dilation confirmed by the
Pound--Rebka experiment, the Gravity Probe~A mission, and the GPS satellite
system.

The Aether density $\rhoa$ and the speed of light $c = \sqrt{\Ka/\rhoa}$
remain constant throughout the inflow.\ This follows from Bernoulli's equation
for the Aether's logarithmic equation of state ($P = \Ka\ln(\rho/\rho_0)$):
along a streamline, $\tfrac{1}{2}\vin^2 + h(\rho) + \Phi = \mathrm{const}$.\
For free-fall inflow ($\tfrac{1}{2}\vin^2 = GM/r$ and $\Phi = -GM/r$,
so $\tfrac{1}{2}\vin^2 + \Phi = 0$), the enthalpy terms cancel exactly,
giving $\rho = \rho_0$ everywhere.\ This result is
equation-of-state independent: for \emph{any} barotropic fluid in free-fall,
the kinetic energy equals the gravitational potential energy, leaving no
energy for compression.

Both velocity and gravitational time dilation now arise from a single
mechanism: motion relative to the local Aether.\ Velocity time dilation
occurs when an object moves through the Aether; gravitational time dilation
occurs when the Aether moves past a static object.\ In both cases, the
fractional time dilation is $v_{\mathrm{rel}}^2/(2c^2)$, where
$v_{\mathrm{rel}}$ is the velocity between the object and the local Aether.\
A freely falling observer---one who co-moves with the inflowing
Aether---experiences $v_{\mathrm{rel}} = 0$ and therefore \emph{no} time
dilation, recovering the equivalence principle from the medium's dynamics.\
$G$ connects directly to the medium: it measures the strength of the inflow
driven by the bind--unbind cycle, expressed in units of the medium's wave
speed ($c^2 = \Ka/\rhoa$).


\subsection{The Event Horizon and Inflow Conservation}

The inflow velocity $\vin(r) = \sqrt{2GM/r}$ (Eq.~\ref{eq:inflow}) is a steady-state
\emph{circulation}, not a one-way accumulation.\ Aether flows inward, participates
in transient binding near the mass, and is returned upon pair annihilation---only
the tiny asymmetry fraction $\epsilon$ escapes as gravitational radiation (see
Section~\ref{sec:steady-state} for the detailed conservation analysis).

This inflow has two further consequences:

\begin{enumerate}[leftmargin=2em]
  \item \textbf{Gravitational lensing.}\ Light propagating through the inflowing
    Aether follows geodesics of the acoustic metric---the effective spacetime
    geometry experienced by waves in a moving medium.\ For an inflow at
    $\vin = \sqrt{2GM/r}$ with uniform wave speed~$c$, the acoustic metric
    takes the Painlev\'e--Gullstrand form, which is exactly equivalent to the
    Schwarzschild geometry.\ The total deflection angle is $4GM/(r_0\,c^2)$,
    matching general relativity (Section~\ref{sec:lensing}).

  \item \textbf{The event horizon.}\ At the Schwarzschild radius $\rs = 2GM/c^2$,
    the inflow velocity equals~$c$.\ Outward-propagating waves can no longer escape,
    producing a sonic horizon analogous to those in supersonic fluid dynamics.\
    Inside this radius, the transient binding cycle still operates locally, but the
    restoring pressure pulses from pair annihilation are swept inward by the
    superluminal flow.
\end{enumerate}


\subsection{Mass--Energy Equivalence (\texorpdfstring{$E = mc^2$}{E = mc²})}

This theory requires that Aether can bind into mass.\ It follows that mass can also unbind---dissolving back into Aether.\ The energy released by this dissolution can be calculated,
and the result connects directly to the transient binding mechanism that drives gravity.
In quantum field theory terms, this dissolution is annihilation---the same process that
ends each virtual particle cycle (and the reverse of the permanent binding produced by
the Dynamic Casimir Effect).

When a mass dissolves into Aether, its volume before dissolution is negligible compared
to its volume as unbound Aether (we are considering the volume of the actual particle
matter, not the empty space within atoms).\ The dissolved Aether must displace the
surrounding ambient Aether, and the total mass of Aether displaced equals the original
mass (imagine a box perfectly filled with identical marbles---forcing one additional marble
in displaces exactly one marble's worth of mass).

This displacement creates a compression wave propagating at the speed of light ($c$), since
light speed is the propagation speed of compression waves in Aether.

The total energy carried by the compression wave follows from the bulk
modulus $\Ka$ and the displaced volume $V = m/\rhoa$:

\begin{itemize}[leftmargin=2em]
  \item \textbf{Elastic energy} from the compression: the energy stored per
    unit volume displaced in a medium of bulk modulus $\Ka$ is
    $\tfrac{1}{2}\Ka$ (at unit strain).\ For volume $V = m/\rhoa$:
    $E_p = \tfrac{1}{2}\Ka\,(m/\rhoa) = \tfrac{1}{2}({\Ka}/{\rhoa})\,m
    = \tfrac{1}{2}mc^2$.
  \item \textbf{Kinetic energy}: for a progressive compression wave in
    a linear medium, the time-averaged kinetic and elastic energies are
    equal (equipartition), so $E_k = E_p = \tfrac{1}{2}mc^2$.
  \item \textbf{Total wave energy}: $E = E_k + E_p = mc^2$.
\end{itemize}

The result depends only on $c^2 = \Ka/\rhoa$: the energy per unit mass
dissolved equals the square of the wave speed.\ This is the standard
energy--mass relation for compression waves in an elastic medium.

In the Aether model, this equation represents the electromagnetic energy released when mass
of negligible pre-dissolution volume dissipates into ambient Aether---and equivalently,
the energy required to compress Aether into mass (the energy cost paid when the DCE
threshold is exceeded and a binding event becomes permanent).\ Each transient binding cycle
in the gravity mechanism momentarily invests $mc^2$ of energy to form a virtual pair and
recovers it upon annihilation.\ The gravitational wave carries only the tiny asymmetry of
this process outward.\ The concept of mass as condensed energy, central to relativity, is
independently predicted by compressible Aether.


\subsection{Relativistic Kinetic Energy}

The preceding sections derived two results purely from Aether mechanics: (1) all
internal processes of an object moving at velocity $v$ through the medium slow by
a factor $\gamma = 1/\sqrt{1 - v^2/c^2}$, from the geometric path-length increase
of internal waves; and (2) $E = mc^2$, the energy stored in a mass at rest.\ These
two results together determine the kinetic energy of a moving object without
invoking any further assumptions.

Consider a particle (a vortex excitation) being accelerated from rest to velocity
$v$ by an applied force $F$.\ The particle's structural integrity is maintained by
its internal circulation---the same wave processes that experience time dilation.\
As the particle moves faster, these internal processes run slower by the factor
$\gamma$, which means the particle's inertial response to the applied force
decreases: its internal dynamics can redistribute momentum only at the dilated
rate.\ The effective momentum of the moving particle is therefore:
%
\begin{equation}
  p = \gamma m v
  \label{eq:rel-momentum}
\end{equation}
%
This is not an additional postulate---it follows directly from the time dilation
already derived.\ A vortex whose internal clock runs $\gamma$ times slower responds
to force $\gamma$ times more sluggishly, producing an effective inertia of
$\gamma m$.

The kinetic energy is the work done to accelerate the particle from rest,
$W = \int_0^v v'\,dp$.\ Evaluating with $p = \gamma mv$ gives:
%
\begin{equation}
  \mathrm{KE} = (\gamma - 1)\,mc^2
\end{equation}
%
At low velocities, $\gamma \approx 1 + v^2/(2c^2)$, recovering the classical
kinetic energy $\tfrac{1}{2}mv^2$.\ As $v \to c$, the energy diverges: the
compression waves that constitute the object's internal structure can no longer
propagate ahead of the object, and infinite energy would be required to reach
the medium's wave speed.

The total energy of a moving particle is thus $E_{\mathrm{total}} = \gamma mc^2$:
rest energy plus kinetic energy, both arising from the mechanics of a vortex
excitation in a compressible medium with finite propagation speed $c$.\ This
completes the triad of $v^2/c^2$ effects---time dilation, magnetic force, and
relativistic energy---all derived from the same Aether physics without recourse to
special-relativistic postulates.


% --------------------------------------------------
\section{Classical Aether Experiments}
\label{sec:classical-experiments}
% ============================================================

The experiments of the 19th and early 20th centuries are widely regarded as having
disproved the existence of Aether.\ However, they disproved only the rigid, static
luminiferous aether of that era---a medium assumed to have no physical effect on the
objects moving through it.\ Dynamic Aether---a medium in which all forces
propagate at speed $c$, producing length contraction, time dilation, and
inflow-dependent gravitational effects---predicts the results of every
classical aether experiment.\ Five key experiments and their explanations
follow.


\subsection{The Michelson--Morley Experiment}

The Michelson--Morley experiment \cite{michelson1887} (1887) measured the round-trip travel time of light along
two perpendicular arms of an interferometer.\ If Earth moves through a static Aether at
velocity $v$, light traveling along the arm parallel to Earth's motion should take longer
than light traveling perpendicular to it, producing an interference fringe shift.

In this theory, the null result is predicted by length contraction.\ An object moving at
velocity $v$ through the Aether is contracted in the direction of motion by the factor
$\sqrt{1 - v^2/c^2}$, because the binding forces holding the interferometer together
propagate at $c$ through the medium and their retarded-field equilibrium contracts by
this factor (Section~\ref{sec:special-relativity}).\ The parallel arm of the
interferometer is shortened by exactly $1/\gamma$, which precisely compensates for the
longer light travel time along that arm.

The round-trip travel time for light along the parallel arm is
$t_{\parallel} = (2L/c)/(1 - v^2/c^2)$, while the perpendicular round-trip
is $t_{\perp} = (2L/c)/\sqrt{1 - v^2/c^2}$ (Section~\ref{sec:special-relativity}).\
These differ by a factor of $\gamma$.\ With the parallel arm contracted to
$L' = L/\gamma$, the contracted parallel time becomes
$t_{\parallel}(L') = (2L'/c)/(1 - v^2/c^2) = (2L/c)/\sqrt{1 - v^2/c^2}
= t_{\perp}$.\ The fringe shift is exactly zero---not because the Aether
does not exist, but because it contracts the apparatus in a way that
conceals its own presence.


\subsection{The Kennedy--Thorndike Experiment}

The Kennedy--Thorndike experiment \cite{kennedy1932} (1932) modified the Michelson--Morley design by using
arms of deliberately different lengths.\ Length contraction alone cannot produce a null
result when the arms are unequal, because the contraction factor cancels only when both
arms are the same length.\ A non-null result would therefore reveal that length contraction
is the sole effect of motion through the medium.

The null result of this experiment is explained by this theory because it predicts both
length contraction and time dilation.\ The combination of spatial contraction (shortening
the parallel arm) and temporal dilation (slowing the clock that measures the fringe
pattern) together produce a complete null result for any arm length ratio.\ This is the
same combination that special relativity invokes, but here both effects arise from the
physical mechanics of moving through a medium in which all forces propagate at
finite speed $c$.


\subsection{The Sagnac Effect}

The Sagnac effect \cite{sagnac1913} (1913) is the observed fringe shift when light travels around a rotating
loop---light traveling with the rotation arrives later than light traveling against it.
Unlike the Michelson--Morley null result, the Sagnac effect is a positive detection of
motion.

This result is more naturally explained by a medium theory than by special relativity.
If light propagates at speed $c$ relative to the Aether (not relative to the apparatus),
then in a rotating interferometer the co-rotating beam must travel a longer path (the
far mirror has moved away) while the counter-rotating beam travels a shorter path (the
near mirror has moved closer).\ The path difference produces a fringe shift proportional
to the angular velocity and the enclosed area:
%
\begin{equation}
  \Delta t = \frac{4A\omega}{c^2}
\end{equation}
%
where $A$ is the enclosed area and $\omega$ is the angular velocity.\ This is exactly the
observed result.\ In the Aether model, the explanation is straightforward: light moves at $c$
relative to the medium, and the rotating apparatus moves relative to both.\ The Sagnac
effect is direct evidence that light speed is determined by a local reference---precisely
what a medium theory predicts.

Modern ring laser gyroscopes and fiber optic gyroscopes exploit this effect for navigation,
confirming the Sagnac fringe shift to extraordinary precision.


\subsection{Stellar Aberration}

Stellar aberration---the apparent displacement of star positions due to Earth's orbital
velocity---was first observed by Bradley in 1728 \cite{bradley1728}.\ A telescope must be tilted slightly in
the direction of Earth's motion to center a star's image, by an angle of approximately
20.5 arc-seconds.

In this theory, the explanation is direct: Earth moves through the Aether at orbital
velocity $v$, while starlight propagates at $c$ through the same medium.\ The telescope must
be tilted by angle $\theta$ where $\tan\theta = v/c$ to compensate for Earth's motion during
the time light traverses the telescope tube.

This result historically posed a dilemma for 19th-century aether theories.\ If the Aether
were dragged along with Earth, aberration should not occur (the medium and telescope would
move together).\ If the Aether were stationary, the Michelson--Morley experiment should
have detected Earth's motion.\ Dynamic Aether resolves this dilemma: Earth moves freely
through the Aether (producing aberration), while length contraction and time dilation
(physical consequences of that motion) eliminate the Michelson--Morley signal.\ The two
results, long considered contradictory for any aether model, are simultaneously explained.


\subsection{The Fizeau Experiment and the Fresnel Drag Coefficient}

The Fizeau experiment \cite{fizeau1851} (1851) measured the speed of light in moving water.\ If water fully
drags the Aether with it, light in the moving water should travel at $c/n + v$ (where $n$ is
the refractive index and $v$ is the water velocity).\ If water does not drag the Aether at
all, light should travel at $c/n$ regardless of water motion.

Fizeau found neither result.\ Instead, the speed of light in moving water is:
%
\begin{equation}
  c_{\mathrm{observed}} = \frac{c}{n} + v\!\left(1 - \frac{1}{n^2}\right)
\end{equation}
%
The factor $(1 - 1/n^2)$ is the Fresnel drag coefficient \cite{fresnel1818}---partial dragging, where only a
fraction of the medium's velocity adds to the light speed.

This result derives naturally from the Aether framework.\ A material with refractive index
$n$ slows light because it increases the effective Aether density within the material's
volume.\ (This is distinct from gravitational fields, where Bernoulli gives
$\Delta\rhoa = 0$.)\ Since $c = \sqrt{\Ka / \rhoa}$ and the material reduces light speed
by factor $n$,
the effective density within the material is:
%
\begin{equation}
  \rho_{\mathrm{material}} = n^2 \, \rhoa
\end{equation}
%
The excess density above the ambient background is:
%
\begin{equation}
  \Delta\rho = (n^2 - 1)\,\rhoa
\end{equation}
%
When the material moves at velocity $v$, only this excess density moves with it.\ The ambient
Aether density $\rhoa$ remains stationary---it is the background medium.\ The fraction of
the total density that is dragged with the material is:
%
\begin{equation}
  f_{\mathrm{drag}} = \frac{\Delta\rho}{\rho_{\mathrm{material}}} = \frac{(n^2 - 1)\,\rhoa}{n^2\,\rhoa}
    = \frac{n^2 - 1}{n^2} = 1 - \frac{1}{n^2}
\end{equation}
%
This is exactly the Fresnel drag coefficient.\ The partial dragging is not an empirical
correction but a direct consequence of the density structure of Aether within refractive
media.\ The ambient Aether is stationary; only the material's contribution to the local
density moves with the material.

This derivation also explains why the drag coefficient depends only on the refractive
index: it is purely a density ratio, independent of the material's chemical composition,
temperature (except through $n$), or other properties.

This derivation carries a significant implication for the bulk modulus $\Ka$.\ The
relationship $\rho_{\mathrm{material}} = n^2 \rhoa$ was derived by assuming that $\Ka$ remains constant
inside the material---that the material changes only the density, not the stiffness, of
the Aether.\ If the material also altered $\Ka$, then $c_{\mathrm{material}} = \sqrt{K_{\mathrm{material}} /
\rho_{\mathrm{material}}}$ and the Fresnel drag coefficient would depend on both the density ratio and
the stiffness ratio, not solely on the refractive index.\ The experimental success of the
$1 - 1/n^2$ coefficient confirms that $\Ka$ is indeed invariant inside materials---the same
conclusion reached independently from velocity time dilation, but here confirmed across
a vastly wider range of density variation.

The density relationship also means that every refractive index measurement is a
measurement of local Aether density.\ Water ($n = 1.33$) multiplies the ambient Aether
density by 1.77.\ Glass ($n = 1.52$) multiplies it by 2.31.\ Diamond ($n = 2.42$) multiplies
it by 5.86.\ The entire field of optics becomes, in this framework, a body of evidence
about how different materials modify the local Aether density---with refractive index
serving as the density probe.


% ============================================================

% ============================================================
\section{Summary of Theoretical Properties}
% ============================================================

\subsection{Properties of Interferometric Predictions}

\begin{itemize}[leftmargin=2em]
  \item The Michelson--Morley null result is predicted by length contraction: the parallel arm
    contracts by $\sqrt{1 - v^2/c^2}$, exactly compensating for the directional light travel
    time difference.
  \item The Kennedy--Thorndike null result requires both length contraction and time dilation,
    both of which this theory independently predicts from motion through the medium.
  \item The Sagnac effect is a direct prediction of light propagating at $c$ relative to the
    medium: rotating apparatus creates an asymmetric path length, producing a fringe shift
    of $\Delta t = 4A\omega/c^2$.
  \item Stellar aberration is explained by Earth's motion through Aether at orbital velocity,
    requiring telescope tilt of $v/c$---approximately 20.5 arc-seconds.
  \item The Fresnel drag coefficient ($1 - 1/n^2$) derives from the density structure of Aether
    in refractive media: material with refractive index $n$ creates effective density
    $n^2 \rhoa$, of which only the excess $(n^2 - 1)\rhoa$ moves with the material while
    ambient $\rhoa$ remains stationary.
  \item The Fresnel drag derivation independently confirms that $\Ka$ (bulk modulus) is invariant
    inside materials: if $\Ka$ changed with density, the drag coefficient would deviate from
    $1 - 1/n^2$.\ This extends the linearity confirmation from density perturbations of order
    $v^2/c^2$ to density ratios up to 6:1.
  \item Every refractive index measurement is a measurement of local Aether density: a material
    with refractive index $n$ multiplies the ambient density by $n^2$.
  \item The gravitational lensing bending angle $4GM/(r_0 c^2)$ arises entirely from the
    acoustic metric of the inflowing Aether: transverse advection by the flow
    ($2GM/(r_0 c^2)$) and time-delay asymmetry from the flow ($2GM/(r_0 c^2)$),
    matching general relativity without spacetime curvature.
  \item The classical aether experiments, once considered fatal to aether theories, are fully
    explained by a medium in which all forces propagate at $c$, producing length contraction
    and time dilation as physical consequences of motion.
\end{itemize}




% ============================================================
\section{Experimental Support and Connections to Established Physics}
% ============================================================

The following experimentally confirmed phenomena support or align with the aspects of the theory presented in this paper:

\begin{enumerate}[leftmargin=2em]
  \item \textbf{Relativistic $\mathbf{E}$/$\mathbf{B}$ field mixing.} In special relativity, a pure electric field in one
    reference frame appears as a combination of electric and magnetic fields in a frame
    moving relative to the first.\ This is a direct consequence of the Aether model: a
    stationary charge has a purely radial flow vector ($\mathbf{E}$ only), but in a frame where
    the charge moves, that flow vector is tilted by the motion, developing angular
    structure ($\mathbf{B}$ appears).

  \item \textbf{Precision velocity time dilation.} The Lorentz factor $\gamma = 1/\sqrt{1 - v^2/c^2}$ has
    been confirmed across velocities from $10^{-6}\,c$ (GPS \cite{ashby2003}, M\"ossbauer \cite{kundig1963}) to $0.9994c$ (CERN muon
    storage ring), with no higher-order corrections ($v^4/c^4$ terms) detected to $10^{-8}$
    precision.\ In the Aether model, this confirms that the medium is perfectly linear in
    its pressure--density response.\ The Weber--Kohlrausch experiment \cite{weber1856} independently confirmed that the $c$ in the
    electromagnetic force ratio is the same $c$ as the speed of light.

  \item \textbf{The Michelson--Morley null result.} The Michelson--Morley experiment \cite{michelson1887} found no directional dependence
    of light speed despite Earth's orbital velocity through space.\ This theory predicts
    exactly this result: length contraction of the parallel interferometer arm by the
    factor $\sqrt{1 - v^2/c^2}$ precisely compensates for the upstream/downstream travel
    time difference, producing zero fringe shift.

  \item \textbf{The Kennedy--Thorndike null result.} The Kennedy--Thorndike experiment \cite{kennedy1932} used unequal arm lengths,
    where length contraction alone cannot produce a null result.\ This theory predicts
    the null result because it provides both length contraction and time dilation---the
    same combination that special relativity requires, but derived from physical motion
    through a medium in which all forces propagate at $c$.

  \item \textbf{The Sagnac effect.} Observed in 1913 \cite{sagnac1913} and confirmed to high precision in modern ring
    laser gyroscopes, the fringe shift in a rotating interferometer
    ($\Delta t = 4A\omega/c^2$) is a direct prediction of light propagating at $c$ relative
    to the medium rather than relative to the apparatus.

  \item \textbf{Stellar aberration.} Bradley's 1728 observation \cite{bradley1728} of 20.5 arc-second stellar displacement
    due to Earth's orbital velocity is directly predicted by Earth's motion through Aether.
    This theory simultaneously explains why aberration occurs (Earth moves through the
    medium) and why the Michelson--Morley experiment is null (length contraction and time
    dilation cancel the signal), resolving a long-standing dilemma for aether models.

  \item \textbf{The Fresnel drag coefficient.} Fizeau's 1851 measurement \cite{fizeau1851} of light speed in moving water
    found partial dragging by the factor $(1 - 1/n^2)$.\ This coefficient derives directly
    from the Aether density framework: material with refractive index $n$ creates effective
    density $n^2 \rhoa$, of which only the excess $(n^2 - 1)\rhoa$ moves with the
    material.

  \item \textbf{$\Ka$ invariance across refractive media.} The Fresnel drag coefficient ($1 - 1/n^2$) is
    exact only if the bulk modulus $\Ka$ remains constant inside materials---if $\Ka$ varied
    with density, the coefficient would acquire additional terms.\ The experimental exactness
    of $1 - 1/n^2$ across materials with widely differing refractive indices confirms $\Ka$
    invariance at density ratios up to $n^2 \approx 6$ (for diamond), extending the linearity
    confirmation by three orders of magnitude beyond the velocity-based experiments.

\end{enumerate}



% ============================================================
\section{Open Questions and Testable Predictions}
% ============================================================

The following areas relevant to this paper remain underdeveloped and represent opportunities for further refinement or falsification:

\begin{enumerate}[leftmargin=2em]
  \item \textbf{Extending the Fizeau framework to dispersive media.} The derivation of the Fresnel
    drag coefficient from $\rho_{\mathrm{material}} = n^2 \rhoa$ assumes a single refractive index.
    In dispersive media, $n$ varies with wavelength, and Lorentz \cite{lorentz1887} showed that the drag
    coefficient acquires a dispersion correction:
    $f = 1 - 1/n^2 - (\lambda/n)(dn/d\lambda)$.\ Deriving this correction from the Aether
    density model would strengthen the connection between refractive index and
    Aether density.

\end{enumerate}



% ============================================================
\section{Mathematical Summary}
% ============================================================

This section collects the key equations relevant to the topics covered in this paper.
\noindent\textit{For the complete mathematical summary, see Paper~1: General Overview \cite{otto-overview}.}

\subsection{Fundamental Aether Parameters}
% --------------------------------------------------

The theory posits four microscopic parameters from which all macroscopic
phenomena derive:
%
\begin{center}
\begin{tabular}{@{}llll@{}}
\hline
\textbf{Symbol} & \textbf{Name} & \textbf{Constrained range} & \textbf{Primary constraint} \\
\hline
$\rhoa$  & Ambient density     & $10^{11}$--$10^{12}$~kg/m$^3$         & Electron vortex energy \\
$\Ka$    & Bulk modulus         & $10^{28}$--$10^{29}$~J/m$^3$          & $\Ka = \rhoa\,c^2$ \\
$m_A$    & Aether quantum mass  & $500$--$850$~MeV/$c^2$                & Flux tube width / glueball mass \\
$\xi$    & Healing length       & $0.17$--$0.28$~fm                     & Bridge diameter (lattice QCD) \\
\hline
\end{tabular}
\end{center}

These four are not independent.\ Two defining relations connect them:
%
\begin{align}
  c &= \sqrt{\Ka / \rhoa}
    \label{eq:sum-speed-of-light} \\
  \xi &= \frac{\hbar}{m_A\,c\,\sqrt{2}}
    \label{eq:sum-healing-length}
\end{align}
%
leaving two free Aether parameters (e.g., $\rhoa$ and $m_A$) from which all
others follow.\ (A third parameter, $\epsilon$, enters through gravity; see below.)


% --------------------------------------------------
\subsection{Equation of State}
% --------------------------------------------------

The Aether obeys a logarithmic equation of state, the unique form that produces
a density-independent bulk modulus (Section~\ref{sec:qss}):
%
\begin{equation}
  P = \Ka \ln\!\Bigl(\frac{\rho}{\rho_0}\Bigr) + P_0
  \label{eq:sum-eos}
\end{equation}
%
Key consequence: the speed of sound $c_s = \sqrt{\Ka/\rho}$ varies with
density.\ At ambient density, $c_s = c$.\ Note: in gravitational fields,
Bernoulli's equation for free-fall inflow gives $\Delta\rho/\rhoa = 0$
exactly, so $c$ does not vary near gravitating masses; the density
dependence is relevant for material media and for the two-fluid
cosmological structure.\ This equation of state guarantees that $\Ka$ is
constant under compression, which is experimentally confirmed by:
%
\begin{itemize}[leftmargin=2em]
  \item Velocity time dilation measurements spanning $10^{-6}\,c$ to $0.9994\,c$
    (precision $\sim 10^{-8}$).
  \item The Fresnel drag coefficient $1 - 1/n^2$, exact across materials with
    density ratios up to $\sim$6:1.
\end{itemize}


% --------------------------------------------------
\subsection{Special Relativity}
% --------------------------------------------------

All relativistic effects arise from motion through the medium.

\paragraph{Lorentz factor.}
%
\begin{equation}
  \gamma = \frac{1}{\sqrt{1 - v^2/c^2}}
  \label{eq:sum-lorentz}
\end{equation}
%
Derived geometrically from the ratio of helical to straight path lengths
in the medium.

\paragraph{Mass--energy equivalence.}
%
\begin{equation}
  E = m\,c^2
  \label{eq:sum-emc2}
\end{equation}
%
Mass is bound Aether; $c^2 = \Ka/\rhoa$ is the medium's elastic energy per
unit density.\ When mass dissolves, it releases a compression wave with this
energy.

\paragraph{Velocity time dilation.}
A moving vortex must traverse a longer helical path through the medium.\ All
internal processes slow by $\gamma$:
%
\begin{equation}
  t_{\mathrm{moving}} = \gamma\,t_{\mathrm{rest}}
  \label{eq:sum-vel-time-dilation}
\end{equation}

\paragraph{Length contraction.}
All binding forces propagate at $c$; the shifted retarded-force equilibrium
contracts bound systems along the direction of motion:
%
\begin{equation}
  L_{\mathrm{moving}} = \frac{L_{\mathrm{rest}}}{\gamma}
  \label{eq:sum-length-contraction}
\end{equation}

\paragraph{Fresnel drag coefficient.}
%
\begin{equation}
  f = 1 - \frac{1}{n^2}
  \label{eq:sum-fresnel}
\end{equation}
%
where $n$ is the refractive index.\ A material with index $n$ has local Aether
density $n^2\,\rhoa$; only the excess $(n^2 - 1)\rhoa$ moves with the material.


% --------------------------------------------------
\subsection{Linearity Constraint}
% --------------------------------------------------

The theory requires that all density perturbations from known particle
interactions are small compared to $\rhoa$.\ The largest perturbation from
the electron's own gravitational field occurs at its Compton wavelength,
where the escape velocity is $v_{\mathrm{esc}} = \sqrt{2Gm_e/\lambdabar_C}$
and the corresponding density ratio is $v_{\mathrm{esc}}^2/c^2$:
%
\begin{equation}
  \frac{\delta\rho}{\rhoa} = \frac{v_{\mathrm{esc}}^2}{c^2}
    = \frac{2\,G\,m_e}{\lambdabar_C\,c^2}
    \approx 3.5 \times 10^{-45} \ll 1
  \label{eq:sum-linearity}
\end{equation}
%
This extraordinary smallness confirms the Lorentz factor to 45 decimal places
and guarantees that the linear approximation underlying special relativity is
valid for all accessible energies.\ Deviations from exact Lorentz symmetry
are predicted only at compressions approaching $\delta\rho/\rhoa \sim 1$,
far beyond current experimental reach.


% --------------------------------------------------

\section*{Acknowledgments}
\addcontentsline{toc}{section}{Acknowledgments}

The author acknowledges the use of Claude (Anthropic) for \LaTeX{} formatting
and editorial assistance.

\bigskip


% ============================================================
% BIBLIOGRAPHY
% ============================================================

\begin{thebibliography}{99}

\bibitem{weber1856}
Weber, W.\ and Kohlrausch, R.,
``Ueber die Elektricit\"{a}tsmenge, welche bei galvanischen Str\"{o}men durch den Querschnitt der Kette fliesst,''
\textit{Annalen der Physik und Chemie}, \textbf{99}, 10--25 (1856).


\bibitem{ashby2003}
Ashby, N.,
``Relativity in the Global Positioning System,''
\textit{Living Reviews in Relativity}, \textbf{6}, 1 (2003).


\bibitem{kundig1963}
K\"{u}ndig, W.,
``Measurement of the Transverse Doppler Effect in an Accelerated System,''
\textit{Physical Review}, \textbf{129}, 2371--2375 (1963).


\bibitem{fresnel1818}
Fresnel, A.,
``Lettre d'Augustin Fresnel \`{a} Fran\c{c}ois Arago, sur l'influence du mouvement terrestre dans quelques ph\'{e}nom\`{e}nes d'optique,''
\textit{Annales de Chimie et de Physique}, \textbf{9}, 57--66 (1818).


\bibitem{lorentz1887}
Lorentz, H.~A.,
``De l'influence du mouvement de la terre sur les ph\'{e}nom\`{e}nes lumineux,''
\textit{Archives n\'{e}erlandaises des sciences exactes et naturelles}, \textbf{21}, 103--176 (1887).


\bibitem{lorentz1904}
Lorentz, H.~A.,
``Electromagnetic phenomena in a system moving with any velocity smaller than that of light,''
\textit{Proceedings of the Royal Netherlands Academy of Arts and Sciences}, \textbf{6}, 809--831 (1904).


\bibitem{michelson1887}
Michelson, A.~A.\ and Morley, E.~W.,
``On the Relative Motion of the Earth and the Luminiferous Ether,''
\textit{American Journal of Science}, \textbf{34}, 333--345 (1887).


\bibitem{kennedy1932}
Kennedy, R.~J.\ and Thorndike, E.~M.,
``Experimental Establishment of the Relativity of Time,''
\textit{Physical Review}, \textbf{42}, 400--418 (1932).


\bibitem{sagnac1913}
Sagnac, G.,
``L'\'{e}ther lumineux d\'{e}montr\'{e} par l'effet du vent relatif d'\'{e}ther dans un interf\'{e}rom\`{e}tre en rotation uniforme,''
\textit{Comptes Rendus}, \textbf{157}, 708--710 (1913).


\bibitem{bradley1728}
Bradley, J.,
``A Letter from the Reverend Mr.\ James Bradley \ldots\ Giving an Account of a New Discovered Motion of the Fix'd Stars,''
\textit{Phil.\ Trans.\ Royal Society}, \textbf{35}, 637--661 (1728).


\bibitem{fizeau1851}
Fizeau, H.,
``Sur les hypoth\`{e}ses relatives \`{a} l'\'{e}ther lumineux,''
\textit{Comptes Rendus}, \textbf{33}, 349--355 (1851).



\bibitem{otto-overview}
Otto, E.,
``Unifying Theory of Dynamic Aether Mechanics: General Overview,''
(2026).


\bibitem{otto-quantum}
Otto, E.,
``Unifying Theory of Dynamic Aether Mechanics: Quantum Mechanics,''
(2026).
\end{thebibliography}

\end{document}